% language: swedish
% figure: fig1 0.12 0.75 0.245 0.858
Massa i volymen $V$
\[ \text{Massa i } V = \iiint_V \rho\,dV \]
\begin{align*}
\text{Massflöde ut ur } V &= \oiint_S \rho\,\vb{u} \cdot \vb{n}\,dS, \\
\text{in i } V &= -\oiint_S \rho\,\vb{u} \cdot \vb{n}\,dS
\end{align*}
\[ \frac{d}{dt} \iiint \rho\,dV = -\oiint \rho\,\vb{u} \cdot \vb{n}\,dS \quad \left\{ \begin{array}{l} \text{byt ordning på} \\ \text{integral och derivata} \end{array} \right\} \]
\begin{align*}
&\Rightarrow \iiint_V \frac{d}{dt}\rho\,dV = -\iiint_V \nabla \cdot (\vb{u}\rho)\,dV \\
&\iiint_V \frac{d\rho}{dt} + \nabla \cdot (\vb{u}\rho)\,dV = 0 \Rightarrow \frac{d\rho}{dt} + \nabla \cdot (\rho\vb{u}) = 0
\end{align*}

Om $\rho(\vb{r}, t)$ konstant i tiden:\\
$\Rightarrow \nabla \cdot (\vb{u}\rho) = 0$

\subsubsection*{Källor}
Integral över sluten yta $S$.
\[ \oiint_S \vb{u} \cdot \vb{n}\,dS = q \]
Om $q \neq 0$ fältet har en källa innanför $S$\\
\phantom{Om} $q = 0$ fältet är källfritt
\[ q = \iiint_V \underbrace{\nabla \cdot \vb{u}}_{\rho \text{ \unsure{rymdkälltäthet}, divergens av } \vb{u}} dV \]

\begin{minipage}{0.22\linewidth}
\notefigure[0.9]{fig1}{}
\end{minipage}
\begin{minipage}{0.76\linewidth}
Låt volymen $V$ minska mot $Q$
\[ \lim_{V \to Q} \oiint_S \vb{u} \cdot \vb{n}\,dS = q \]
\end{minipage}

Om gränsvärdet existerar och har $q \neq 0$ så har fältet en punktkälla.
